\section*{Problem 4}

\textbf{(14 points total)} As you may have seen in Lab 0, the function generators in the
laboratory can be simply modeled by a voltage source $V_{\mathrm{gen}}$ with amplitude
$V_m$ in series with a $50\,\Omega$ source resistance; this is illustrated in the circuit
on the right. By default, the function generator ``assumes'' it will be connected to a
matched load; that is, a load $R_L$ whose resistance matches that of the source---namely
$50\,\Omega$.

% Insert draw.io figure here.
%
% \begin{figure}[ht]
%     \centering
%     \includegraphics[width=0.35\textwidth]{figures/problem4_function_generator.pdf}
%     \caption{Function-generator model with a $50\,\Omega$ source resistance.}
%     \label{fig:problem4-function-generator}
% \end{figure}

\begin{enumerate}
\item[(a)] \textbf{In a matched-load scenario, if the output (load) voltage is configured
to have an amplitude of $10\,\mathrm{V}$, what is the actual value of $V_m$ set by the
function generator? (2 points)}

Because the function generator is expecting a matching load of $50\,\Omega$, it expects
to deliver what is configured, $10\,\mathrm{V}$, at the load. Since
\[
V_L = V_{\mathrm{gen}}\frac{R_L}{50\,\Omega + R_L},
\]
we have
\[
V_{\mathrm{gen}}
=V_L\frac{50\,\Omega+R_L}{R_L}
=(10\,\mathrm{V})\frac{50\,\Omega+50\,\Omega}{50\,\Omega}
=20\,\mathrm{V}.
\]

\item[(b)] \textbf{What percentage of the power in the circuit is dissipated in the load
in this scenario? (2 points)}

The power dissipated in the load is
\[
P_L=\frac{V_L^2}{R_L}
=\frac{(10\,\mathrm{V})^2}{50\,\Omega}
=2\,\mathrm{W}.
\]
The total power dissipated is
\[
P_{\mathrm{total}}
=\frac{V_{\mathrm{gen}}^2}{50\,\Omega+R_L}
=\frac{(20\,\mathrm{V})^2}{100\,\Omega}
=4\,\mathrm{W},
\]
so the percentage dissipated in the load is
\[
\frac{P_L}{P_{\mathrm{total}}}
=\frac{2\,\mathrm{W}}{4\,\mathrm{W}}
=0.5=50\%.
\]

\item[(c)] \textbf{Now assume the load $R_L$ is not $50\,\Omega$ but a
$1\,\mathrm{M}\Omega$ oscilloscope channel input. If the function generator remains
configured as in part (a), what is the amplitude of the voltage that will be displayed
on the oscilloscope? (2 points)}

The source voltage remains $V_{\mathrm{gen}}=20\,\mathrm{V}$, but now
\[
V_L
=V_{\mathrm{gen}}\frac{1\,\mathrm{M}\Omega}
{50\,\Omega+1\,\mathrm{M}\Omega}
\approx V_{\mathrm{gen}}
=20\,\mathrm{V}.
\]

\item[(d)] \textbf{How should the function-generator voltage be adjusted to provide the
expected voltage across the load? (2 points)}

We would need to drop the function-generator setting to $5\,\mathrm{V}$, so that
\[
V_{\mathrm{gen}}=V_L=10\,\mathrm{V}.
\]

\item[(e)] \textbf{What percentage of the power in the circuit is now dissipated in the
load? (2 points)}

The power dissipated in the load is
\[
P_L=\frac{V_L^2}{R_L}
=\frac{(10\,\mathrm{V})^2}{1\,\mathrm{M}\Omega}
=0.1\,\mathrm{mW}.
\]
The total power dissipated is
\[
P_{\mathrm{total}}
=\frac{V_{\mathrm{gen}}^2}{50\,\Omega+R_L}
=\frac{(10\,\mathrm{V})^2}{50\,\Omega+1\,\mathrm{M}\Omega}
\approx 0.1\,\mathrm{mW},
\]
so the percentage dissipated in the load is approximately
\[
\frac{0.1\,\mathrm{mW}}{0.1\,\mathrm{mW}}
=1.0=100\%.
\]

\item[(f)] \textbf{What is the approximate source resistance of a $20{:}1$
oscilloscope probe? (2 points)}

We require the voltage provided to the oscilloscope to be $1/20$ of the actual measured
voltage:
\[
\frac{V_{\mathrm{scope}}}{V_{\mathrm{probe}}}
=\frac{R_{\mathrm{scope}}}{R_{\mathrm{probe}}+R_{\mathrm{scope}}}
=\frac{1}{20}.
\]
Since the oscilloscope has a $1\,\mathrm{M}\Omega$ input,
\[
\frac{1\,\mathrm{M}\Omega}{R_{\mathrm{probe}}+1\,\mathrm{M}\Omega}
=\frac{1}{20},
\]
and therefore
\[
R_{\mathrm{probe}}=19\,\mathrm{M}\Omega.
\]

\item[(g)] \textbf{Now imagine you took such a probe attached (and properly configured)
to a second identical oscilloscope and measured the voltage of the load $R_L$ coming into
the first oscilloscope still configured as in part (d). What is the voltage level seen by
the second oscilloscope? What is the voltage level displayed by the second oscilloscope?
(2 points)}

The voltage seen by the second oscilloscope is $1/20$ of the actual voltage, which was
determined to be $10\,\mathrm{V}$. Thus,
\[
V_{\mathrm{seen}}=0.5\,\mathrm{V}.
\]
Because the second oscilloscope is properly configured with the $20{:}1$ probe attenuation
factor, it compensates for this factor of 20 and displays
\[
V_{\mathrm{displayed}}=10\,\mathrm{V}.
\]

\end{enumerate}
